%---------------------------------------------------------------------
\subsubsection*{Déroulement de carrière}
%---------------------------------------------------------------------

Ma carrière jusqu'à présent se découpe en trois grandes parties. À l'issue de ma thèse, j'ai d'abord été maître de conférences puis professeur (PR2 en 2003) à \APT de 1994 à 2007. En 2008, je suis devenu directeur de recherches (DR2, et DR1 en 2016) à l'INRAE jusqu'en 2021, date à laquelle j'ai rejoint \SU en tant que professeur (PR1) sur un poste 46.4 intitulé ``Statistique et Biologie''. 

De 1994 à 2021, j'ai été rattaché à l'UMR de Mathématiques et Informatique Appliquées (MIA : \APT/INRAE/Paris-Saclay), que j'ai dirigée de 2004 à 2013. 
J'ai effectué un détachement à l'INRAE de Jouy-en-Josas de 1999 à 2001, pendant lequel j'ai participé à la création de l'unité INRAE Mathématiques, Informatique et Génome (MIG). 
J'ai également été mis à disposition du Centre d'Écologie et des Sciences de la Conservation (CESCO) du Muséum National d'Histoire Naturelle en 2020 et 2021. J'ai rejoins le Laboratoire de Probabilités, Statistique et Modélisation (LPSM) de \SU en 2021.

%---------------------------------------------------------------------
\subsubsection*{\'Evolution des activités}
%---------------------------------------------------------------------

%---------------------------------------------------------------------
\paragraph{Enseignement.}
% APT
J'ai commencé ma carrière comme enseignant dans une école d'ingénieur, \APT, dont les étudiants sont pour une très large part issus de classes préparatoires BCPST (biologie, physique, chimie, sciences de la vie et de la Terre). Je me suis donc adressé très tôt à un public d’étudiants en sciences, mais dans des filières non principalement mathématiques. J’ai également eu très tôt à participer puis à concevoir des enseignements dans un cadre interdisciplinaire, en collaboration avec des collègues biologistes, agronomes ou écologues. 

Les statuts de maître de conférences et de professeur du ministère de l’Agriculture (auquel est rattachée \APT) sont calqués sur ceux de l’université. Durant cette période, j'ai donc assuré un service plein de 192h équivalent TD par an. Mes enseignements s’adressaient à tous les niveaux de l’école (1ère et 2ème année --~niveaux L3 et M1~--, année de spécialisation et master 2, école doctorale). Ils visaient à donner aux étudiants une formation statistique généraliste (statistiques descriptive, statistique inférentielle, modèles linéaires, méthodes multivariés), mais certains enseignements étaient plus particulièrement associés à certaines spécialisations (planification expérimentale en agronomie, données censurées en santé humaine, modèles mixtes en sciences animales, modèles linéaires généralisés en écologie). Je donnais également différents cours optionnels en probabilités (modélisation aléatoire, modèles markoviens, processus de naissances et morts). De 2004 à 2013, j'ai donné le cours de modèles linéaires du master 2 de Statistique de l'université d'Orsay.

\bigskip
% INRAE
J’ai poursuivi une partie de mes enseignements lorsque j’ai été recruté à l’INRAE. Notamment, j'ai participé à la conception et à création du master ‘Mathématiques pour les sciences du vivant’ (MathSV) de l’université Paris-Saclay/École polytechnique créé en 2012. L’objectif de ce master est de former de jeunes chercheuses et chercheurs en mathématiques appliquées intéressés par les sciences du vivant, en leur fournissant des bases solides en modélisation déterministe et en modélisation aléatoire, ainsi qu'en optimisation. De 2012 à 2021, j’ai proposé dans ce master un cours sur les modèles à variables latentes orientés vers la biologie, la génomique, puis vers l’évolution et l’écologie. J’ai repris en grande partie ce cours en rejoignant \SU.

Pendant cette même période, j’ai continué à participer à un cours de modélisation statistique que j’avais co-créé pour l’école doctorale ABIES (agriculture, alimentation, biologie, environnement et santé). L’objectif de ce cours, réparti sur deux semaines pleines, est de former des doctorants en sciences du vivant à une démarche complète de modélisation, d'analyse statistique d’un problème issu de leurs recherches. Cette formation résultait du constat qu’il est difficile de former les étudiants à une telle démarche tant que les questions de recherche auxquelles elle doit permettre de répondre ne sont pas les leurs. J'ai cessé de participer à cette formation en rejoignant \SU.

Sur l'ensemble de cette période, j'ai assuré l'équivalent d'environ un quart de service d'enseignement universitaire par an.

\bigskip
% \SU
Je décris mon activité d’enseignement à \SU dans la section \ref{sec:ens}. Ces enseignements se répartissent essentiellement en un cours L1 en Sciences des Données, un cours de L3 de Mathématiques pour les biologistes, d’un cours de M2 sur des modèles statistiques pour l’écologie et une série de cours dans un M2 de biologie évolutive et génétique des populations. Ces trois derniers enseignements bénéficient évidement de mon expérience passée en enseignement et en recherche à l’interface entre mathématiques et sciences du vivant.

%---------------------------------------------------------------------
\paragraph{Recherche.}
J’effectue mes recherches dans le domaine de la statistique et des probabilités appliquées, le plus souvent en lien avec les sciences du vivant. Mon expertise porte notamment sur les modèles à variables latentes, les problèmes de détection de rupture et les modèles de réseaux aléatoires. Ma démarche consiste à proposer des modèles statistiques pertinents pour répondre à une question biologique spécifique en se fondant sur des données particulières et à proposer des méthodes d’inférence présentant des garanties statistiques tout en étant efficaces d’un point de vue computationnel.

\bigskip
La première partie de ma carrière de recherche a été fortement liée à l’émergence des données “omiques” (génomiques, transcriptomiques, protéomiques, métagénomiques, etc.). Mon mémoire d’HDR a ainsi porté sur des modèles probabilistes permettant de caractériser la distribution de motifs d'intérêt le long du génome.

\bigskip
J’ai choisi de me réorienter vers les applications en écologie il y a une dizaine d’années. Cette transition s’est faite en partie à l’occasion de collaborations en écologie microbienne qui utilise des données moléculaires pour caractériser des écosystèmes (en l'occurrence microbiens). Cette ré-orientation était aussi motivée par le renouvellement sans cesse accéléré des technologies de biologie moléculaire, qui tend à rendre rapidement caduques certains développements méthodologiques. De plus, mes thématiques de prédilection ont rapidement trouvé des applications en écologie, comme j’en donne quelques exemples en section \ref{sec:rech}.

\bigskip
Au cours de ces années, j’ai pu constituer un réseau de collaboration riche et varié en mathématiques comme en sciences du vivant. En arrivant au LPSM, j’ai pu nouer de nouvelles collaborations autour de thématiques qui m’occupaient déjà. L’arrivée dans une nouvelle équipe m’a aussi permis de m'intéresser à de nouveaux objets comme les processus Hawkes (dont les applications en sciences du vivant sont légion) ou la statistique robuste. Je décris ces collaboration de section \ref{sec:rech}.

%---------------------------------------------------------------------
\paragraph{Responsabilités collectives et animation de la recherche.}
Je décris à la section \ref{sec:col} certaines des responsabilités collectives que j’ai exercées par le passé, ainsi que des actions d’animation scientifique que j’ai pu mener. Je n’occupe aucune responsabilité collective formelle à \SU, mais je contribue néanmoins à l'organisation collective de l’université. 

J’ai notamment participé à la réforme du cours de mathématiques commun en L1 et à la conception du contenu de la mineure de Sciences des Données proposées en L2 et L3 lors de sa création. Je suis par ailleurs membre du comité de pilotage de la Fondation pour la Science Mathématiques de Paris-centre (FSMP) et du comité exécutif de l’initiative InLife (Sciences aux Interfaces du Vivant) de \SU. 

% Je donne le détail de mes responsabilités collectives à la section \ref{sec:col}. 
% La direction de l’équipe \SG de l’\UMRcourt a été ma première responsabilité collective en matière de recherche, il s’agissait d’une responsabilité essentiellement scientifique dans une équipe d’environ 10 membres travaillant sur des sujets proches. La responsabilité la plus importante est certainement la direction de l’\UMRcourt pendant 9 ans (trois évaluations de l’HCERES). L’unité de taille moyenne comptait alors un peu plus de 20 permanents et de 30 membres au total répartis en deux ou trois équipes selon les périodes et travaillant sur des sujets beaucoup plus variés en informatique, statistique et probabilités appliqués avec des applications en risque sanitaire, environnement ou génomique.  

% Pendant mon long séjour à \APT, j’ai toujours été élu dans une instance telle que le conseil des enseignants, les conseils scientifiques ou le conseil d’administration. J’ai également siégé dans les conseils scientifiques de plusieurs départements de l’INRAE, ainsi qu’au CNU, en section 26. J'ai enfin été membre du comité de pilotage de la fondation mathématique Jacques Hadamard (FMJH) et de l'école doctorale de mathématiques Hadamard (EDMH) lors de la création. Je suis ensuite devenu membre du conseil scientifique de la FMJH. Je n’ai assuré aucune responsabilité collective notable depuis mon arrivée à \SU, mais je suis membre du comité de pilotage de la fondation pour la science mathématique de Paris centre (FSMP) et du comité exécutif de l'intiative \SU InLife.

% Pour ce qui est de l’enseignement, comme indiqué plus haut, j’ai participé à la conception, la création et à la mise en place du master ‘Mathématiques pour les Sciences du Vivant” de l’université Paris-Saclay / École polytechnique et à la création et mise en place du master EvoGEM auquel sont associées toutes les universités parisiennes, ainsi que l’université Paris-Saclay.

%---------------------------------------------------------------------
% \subsubsection*{Formations suivies}
%---------------------------------------------------------------------
% Présentation des formations suivies notamment concernant vos activités pédagogiques
